\documentclass[10pt,a4paper,sans]{moderncv}
\moderncvstyle{classic}
\moderncvcolor{blue}
\usepackage[utf8]{inputenc}
\usepackage[scale=0.85]{geometry}
\usepackage{helvet}

\name{Alexandro}{Disla}
\title{Suivi, Évaluation, Redevabilité et Apprentissage (MEAL)}
\address{Rue Saint-Louis Jeanty \#14, Route de Frère}{Port-au-Prince, Haïti}{}
\phone[mobile]{(+509) 4148-3700}
\email{alexandrodisla@hotmail.com}
\extrainfo{GitHub: \url{https://github.com/AD0791}}

\begin{document}
\makecvtitle

\section{Profil Professionnel}
\cvitem{}{Professionnel MEAL avec \textbf{cinq ans d'expérience terrain} sur des projets multi-sites en Haïti, répartis sur trois secteurs --- santé et VIH à la Caris Foundation International, eau et assainissement chez HANWASH, éducation chez Anseye Pou Ayiti. Mon travail couvre le cycle MEAL complet d'un projet : élaboration des indicateurs, conception des outils de collecte, analyse quantitative et qualitative, et restitution aux équipes et aux bailleurs. Praticien de l'\textbf{ICT4D} (ODK, KoboToolbox, CommCare, mWater) et \textbf{ingénieur de données} en parallèle, j'administre les bases et construis les tableaux de bord dont dépend un système MEAL, au lieu de seulement les spécifier. \textbf{Français et créole haïtien maternels}, anglais courant à l'écrit comme à l'oral. Basé à Port-au-Prince, disponible pour Les Cayes ou Fort-Liberté et pour des déplacements fréquents.}

\section{Compétences Clés}
\cvitem{}{
\begin{itemize}
    \item \textbf{MEAL sur le cycle de projet :} Appui au démarrage, à la mise en œuvre et à la clôture ; cadres logiques, chaînes de résultats, plans MEAL et tableaux de suivi des indicateurs ; élaboration d'indicateurs avec valeurs de référence, cibles, désagrégations et sources de vérification ; conformité aux exigences MEAL du bailleur.
    \item \textbf{ICT4D \& gestion des données :} Conception d'instruments XLSForm avec logiques de saut et contraintes de validation sur ODK, KoboToolbox, CommCare et mWater ; numérisation des flux de collecte, administration des bases relationnelles, contrôle d'accès par rôle et stockage sécurisé ; dédoublonnage des listes de bénéficiaires.
    \item \textbf{Analyse quantitative et qualitative :} Analyse statistique et échantillonnage (économiste appliqué de formation) ; \textbf{R}, SQL et Python pour l'analyse et l'automatisation du rapportage ; \textbf{Power BI}, Looker Studio et Quarto pour la visualisation et le suivi périodique des indicateurs.
    \item \textbf{Qualité des données :} Protocoles DQA établis et supervisés, vérification par remontée aux registres sources, suivi des actions correctives jusqu'à clôture --- et non jusqu'à leur simple consignation.
    \item \textbf{Transfert de connaissances \& partenaires :} Rédaction de guides techniques et de supports de formation, formation d'équipes de terrain et de partenaires aux outils de collecte et à la qualité des données ; encadrement technique de personnel de suivi-évaluation ; coordination avec le MSPP, la DINEPA et le MEF.
\end{itemize}}

\section{Expériences MEAL}

\cventry{Oct 2021 -- Août 2024}{Officier de Suivi \& Évaluation}{Impact Youth Project, Caris Foundation International}{Haïti}{}{
\begin{itemize}
    \item \textbf{Système MEAL multi-sites :} Conduite du cycle de suivi d'un projet de santé et de VIH sur des cliniques multi-sites, avec standardisation des formulaires et des flux de données entre sites et consolidation des remontées dans les délais.
    \item \textbf{Encadrement technique :} Direction technique des agents de collecte sur l'ensemble des sites --- répartition du travail, contrôle qualité, formation aux procédures --- et référent technique du personnel de suivi-évaluation sur les contrôles de cohérence.
    \item \textbf{Assurance qualité (DQA) :} Établissement et supervision des protocoles de Data Quality Assessment sur les données intégrées MySQL, avec remontée des chiffres rapportés jusqu'aux registres sources et suivi des actions correctives jusqu'à clôture.
    \item \textbf{ICT4D :} Mise en place et administration des systèmes de collecte numérique alimentant le rapportage (CommCare, HIVHaiti, MySQL).
    \item \textbf{Rapportage bailleur :} Production et soumission des indicateurs MER destinés à l'USAID via DATIM, la plateforme du PEPFAR bâtie sur DHIS2.
    \item \textbf{Analyse \& restitution :} Rapports d'activité automatisés en Python et SQL, réduisant le temps de traitement manuel de 40 \% ; tableaux de bord programme sur AWS (Quarto) et Power BI pour le suivi périodique des indicateurs.
\end{itemize}}

\cventry{Oct 2023 -- Jan 2026}{Consultant MEAL \& Analyste de Données}{HANWASH}{Haïti / Télétravail}{}{
\begin{itemize}
    \item Raffinement des plans MEAL et des tableaux de suivi des indicateurs alignés sur les normes JMP (OMS/UNICEF) pour des programmes d'eau, d'hygiène et d'assainissement.
    \item Déploiement d'enquêtes géospatiales complexes sur mWater et administration de la base ; tableaux de bord automatisés (mWater, Power BI) connectés aux API de collecte, guides techniques et formation des équipes à leur usage autonome.
\end{itemize}}

\cventry{Jan 2026 -- Mars 2026}{Analyste de Données \& Consultant MEAL}{Anseye Pou Ayiti}{Haïti}{}{
\begin{itemize}
    \item Définition de la stratégie de collecte et élaboration des indicateurs de performance de trois programmes ; codage des instruments XLSForm pour ODK avec logiques de saut et contraintes de validation, administration du système ODK, et formation des agents de terrain.
\end{itemize}}

\cventry{Nov 2015 -- Août 2023}{Économiste --- Direction de Planification Économique et Sociale (DPES)}{Ministère de la Planification et de la Coopération Externe (MPCE)}{Port-au-Prince}{}{Suivi de l'exécution du Plan Triennal d'Investissement 2014--2016 : indicateurs de projets publics sur l'ensemble des secteurs et des départements géographiques ; unité statistique et modélisation, avec participation à la construction du << Modèle de Planification du Développement >>, outil de prévision, de simulation et d'analyse d'impact des politiques publiques.}

\section{Expériences en Ingénierie de Données et Systèmes}

\cventry{Févr 2026 -- Présent}{Ingénieur Logiciel (Fullstack)}{Tekkod LLC}{Télétravail}{}{Modernisation d'infrastructure mobile avec fiabilité hors ligne et sécurité des données pour un usage terrain ; tableaux de bord Looker Studio.}

\cventry{Mars 2021 -- Mai 2024}{Développeur Backend \& Ingénieur de Données}{Tekkod LLC}{Télétravail}{}{Administration de bases relationnelles en accès concurrent ; pipelines ETL (Apache Beam) migrant les données MongoDB vers BigQuery ; services REST sécurisés (FastAPI, JWT) avec contrôle d'accès par rôle.}

\cventry{Nov 2015 -- Mars 2021}{Consultant Logiciel \& Spécialiste Intégration de Données}{CassionSoft (Projet UNOPS)}{Haïti}{}{Intégration de données sur un projet sous contrat des Nations Unies, avec exposition aux procédures et standards documentaires onusiens.}

\section{Formation}
\cvitem{2010--2014}{\textbf{Études Supérieures en Économie Appliquée} --- C.T.P.E.A (Centre de Techniques de Planification et d'Économie Appliquée), Port-au-Prince : cycle de quatre ans en économie et statistique appliquées, de niveau licence. \textit{Scolarité complétée ; Diplôme d'Études Supérieures en attente de la soutenance du mémoire de sortie --- attestation officielle jointe au dossier.}}
\cvitem{2009--2010}{Baccalauréat (Mention) --- Institution Saint-Louis de Gonzague, Port-au-Prince}
\cvitem{Autoformation}{FlatIron School (Software Engineering), CodeWithMosh, Udemy}

\section{Compétences Techniques}
\cvitem{ICT4D \& collecte}{\textbf{ODK (Open Data Kit), KoboToolbox, XLSForm}, CommCare, mWater, HIVHaiti, DHIS2/DATIM \textit{(soumission de données)}, protocoles DQA, ITT}
\cvitem{Analyse statistique}{\textbf{R} (RMarkdown/Quarto), SQL, Python (Pandas), modélisation statistique, échantillonnage, analyse qualitative et quantitative}
\cvitem{Bureautique \& viz}{\textbf{MS Excel (avancé), Word, PowerPoint}, GSuite, \textbf{Power BI (expert)}, Looker Studio, Teams et Zoom}
\cvitem{Bases de données}{MySQL, PostgreSQL, MongoDB, SQLite, Azure, AWS, Google Cloud (BigQuery), Docker, Git}
\cvitem{Langues}{\textbf{Français (maternel)}, \textbf{Créole haïtien (maternel)}, \textbf{Anglais (courant, rédaction de rapports)}}

\section{Référence Professionnelle}
\cvitem{}{Davidson Adrien --- ancien Directeur Suivi \& Évaluation, Caris Foundation International --- (+509) 4460-6638}

\end{document}
